\section{Introduction and context}
\label{sec:intro}

This chapter states the objective, the platform, and the scope that survived the year. It closes
with the contributions and a reading map. The rest of the document is then a single justification
chain: physics and measurement first, then constraints, then the design that follows from them.

\subsection{The platform and the objective}
\label{sec:intro-objective}

DASH-01 is a bipedal robot built at REHAssist. It carries two legs, six quasi-direct-drive
actuators, a parallel 4-bar knee, a passive ankle spring and a single-board computer. The hardware
budget was roughly EUR\,5000. \Cref{fig:robot} shows the machine as modelled.

The objective at kickoff was maximum running speed, with the bipedal 100\,m record as the stated
ambition. That record stands at 24.73\,s, or about 4.1\,m/s average, set by Cassie in
2022~\cite{crowley2024cassie}.

The project inherited a partially finished mechanical design. The work reported here therefore
starts at modelling, and at the selection of the non-mechanical subsystems. It then runs through
characterisation, learning and deployment. Five threads organise it:

\begin{enumerate}[leftmargin=*]
  \item finish the robot: actuator, battery, sensor and compute selection;
  \item learn reinforcement learning and the sim-to-real problem on a toy plant, while the robot
        did not yet exist;
  \item model the robot at a fidelity that policies can be trained against;
  \item characterise the physical robot, to bound the sim-to-real risk;
  \item produce a policy that works.
\end{enumerate}

\subsection{What the year changed about the objective}
\label{sec:intro-reframe}

Two facts reframed the objective. Both are stated here rather than conceded late.

\textbf{The record moves faster than one thesis.} New bipedal speed marks now appear at roughly a
monthly cadence, from teams with institutional budgets. A one-person master thesis on a ~CHF\,5000
chassis is not a credible contender on that timescale.

\textbf{The measured plant does not admit record pace.} \Cref{sec:speed-ceiling} bounds this
mechanism at about 2.7\,m/s in burst, and about 1.3\,m/s thermally sustained. The bound comes from
first principles, with no controller in the calculation. A 100\,m dash at record pace is about
25\,s of continuous running, so the sustained figure is the one that decides a record. The gap is a
factor of three, and it is mechanical.

The objective is therefore restated as a measured demonstrator plus a reusable method. The
deliverable is the chain from CAD to a trained policy, with every load-bearing number measured
rather than assumed. It includes an explicit account of the distance to the original target. The
thesis title still names the record, and \cref{sec:future} names the design changes that chasing it
would need.

\textbf{The problem this addresses.} Reinforcement learning has produced legged-locomotion results
on research platforms. Reproducing them on a new, custom, imperfectly known robot is a different
problem. The plant parameters are uncertain. The actuators hide undocumented internal control
loops. The sensors are commodity parts, and the computer is a single board. The gap between
``reinforcement learning walks in simulation'' and ``reinforcement learning walks on this robot'' is
dominated by modelling fidelity, actuator bandwidth, calibration and instrumentation. This thesis
measures each of those.

\textbf{Sensing scope.} Proprioception and one base inertial measurement unit are the only sensors.
Foot impact sensors were dropped for time rather than on evidence. The hardware provision for them
remains in place, as four analogue inputs on the sensor hat. The literature prior that foot feedback
helps convergence is therefore untested here.

\subsection{Contributions}
\label{sec:intro-contrib}

\begin{itemize}[leftmargin=*]
  \item \textbf{A CAD-to-MJCF pipeline for a leg that closes a kinematic loop}, with the
        loop-closure treatment URDF cannot express. It ships with validation tooling: analytic
        Jacobians checked against finite differences, and loop closure checked against the measured
        backdriven workspace at 2789 of 2789 poses (\cref{sec:twin}).
  \item \textbf{A measured plant} (\cref{tab:plant-id}). It covers weighed segment masses, a torque
        constant identified by a known-mass method that breaks the degeneracy between $K_t$ and
        gravity, the Bode response of the drive's internal position loop, inertial-sensor noise
        with provenance labels, loop timing and the CAN protocol. The identification also states
        what it could \emph{not} resolve.
  \item \textbf{A first-principles capability analysis} of running speed on that plant, validated
        against MuJoCo. Its mass-distribution sensitivity ranks the design levers
        (\cref{sec:speed-ceiling}).
  \item \textbf{An action-space and curriculum design} for the learned controller. The generator is
        a Fourier gait basis, re-parameterised every control step and mirrored across legs. It is
        wrapped in hard-wired and learned reflexes and a per-step residual channel, and trained
        under a base-degree-of-freedom release curriculum
        (\cref{sec:rl-cpg,sec:rl-curriculum}).
  \item \textbf{Two controlled comparisons that isolate the plant.} A generator A/B against an
        explicit-oscillator central pattern generator, on identical plants and rewards
        (\cref{sec:cpg-ab}). And a 21-gain scripted classical controller, which beats the learned
        reference on the constrained plant and fails identically on the free one
        (\cref{sec:scripted-baseline}).
  \item \textbf{A failure catalogue with measured diagnoses.} It covers reward exploits, six
        simulation artefacts that each produced a clean converged wrong result, four evaluation
        bugs, and the pitch-balance wall traced to plant statics (\cref{sec:negative}).
  \item \textbf{A deployment stack on the robot.} It is a 200\,Hz daemon on stock Linux, a web
        bring-up instrument that carried most of the measurement campaign, and a three-tier flight
        recorder with a pre-move guard. The recorder was built after a hardware failure, and
        validated by reproducing it (\cref{sec:deploy}).
  \item \textbf{Two published component databases with figure-of-merit tooling}, covering 140
        commercial actuators and roughly 430 battery cells. Both audited this robot's own hardware
        choices after the fact (\cref{sec:hw-selection}, \cref{sec:app-db}).
\end{itemize}

\subsection{Reading map}
\label{sec:intro-map}

\Cref{sec:bg} fixes the vocabulary, reads the field as six answers to one question, and sorts those
answers on seven analytical axes. It closes with the gap this work sits in, and three hypotheses.
\Cref{sec:robot} builds the robot and its digital twin, and replaces every assumed plant parameter
with a measured one. \Cref{sec:rl} develops the controller: five action spaces, one curriculum, and
the measured verdict on each. \Cref{sec:results} reports what the robot and the policies do.
\Cref{sec:discussion} measures the distance to the objective, and ranks the levers that close it.
\Cref{sec:conclusion} recaps. The appendix holds the logbook, the pipeline procedure, and the
database and configuration detail.

One motif recurs across the chapters, and it is worth naming in advance. The binding constraint
moves from the algorithm to the plant. \Cref{sec:bg} introduces it as a demand each control family
places on the hardware. \Cref{sec:robot} demonstrates it, because every measurement moved
pessimistically. \Cref{sec:rl} confirms it, because a learned controller and a scripted one hit the
same wall at the same time. \Cref{sec:discussion} quantifies it.
